% ===========================================================================
%  sviluppatore/06-dashboard/08-overview.tex
% ===========================================================================
\section{La Overview}\label{sez:overview-pagina}

La pagina che si apre per prima. \file{templates/home.html},
\file{static/js/home.js}, \file{static/css/home.css}; la rotta è
\simbolo{index} in \file{dashboard/app.py}. Mostra il valore fuori campione di
un progetto con SPY e il drawdown, le metriche scelte, i pesi, gli anni, i
passi della guida e la tabella di tutti gli stream. I numeri arrivano tutti
dalle API che usa anche il resto dell'interfaccia: la pagina non calcola
niente che il server non sappia già, salvo la finestra del periodo e
qualche somma per la tabella dei pesi.

\subsection{Le sezioni sulla pagina}

Nessun pannello contiene le sezioni: le due colonne sono figlie dirette di
\cod{\#ov-root}, che è la griglia, e né la griglia né le colonne né le sezioni
hanno un fondo proprio. Le separano solo bordi da un pixel: \cod{border-left}
sulla colonna di destra, \cod{border-top} sulla seconda sezione di ogni
colonna. Due colonne, e ciascuna si divide per conto suo:

\begin{center}
\begin{tabular}{@{}L{34mm}L{104mm}@{}}
\toprule
\textbf{Sezione} & \textbf{Che cosa} \\
\midrule
\rigaTenue
\cod{.ovx-port} & Il progetto: nome e menu, giudizio di robustezza, periodo,
  valore e variazione, legenda, grafico. Prende \num{1,32} parti
  dell'altezza della colonna. \\
\cod{.ovx-streams} & La tabella degli stream, che scorre da sola. Una parte. \\
\rigaTenue
\cod{.ovx-side} & Le metriche: alta quanto serve. \\
\cod{.ovx-more} & Le tre schede (pesi, anni, passi): il resto della colonna,
  e scorre. \\
\bottomrule
\end{tabular}
\end{center}

Le linee arrivano ai bordi della pagina, il testo resta allineato al titolo.
\cod{\#ov-root} esce dal margine del canvas con un margine negativo
(\cod{--ovx-gut}: 36 pixel, 20 fino a \num{1024}, gli stessi di
\file{quantsys-modern.css}) e ogni sezione lo rimette come spazio interno con
due variabili, \cod{--ovx-pl} a sinistra e \cod{--ovx-pr} a destra; verso la
linea fra le colonne lo spazio è \cod{--ovx-in}, 24 pixel. Così le righe della
tabella e la loro evidenziazione vanno da bordo a bordo, mentre il nome dello
stream sta sotto il titolo della pagina. Il grafico tiene a destra
\cod{--ovx-pr} meno i 12 pixel che l'SVG lascia già dopo le etichette della
scala: le etichette finiscono dove finisce la scelta del periodo.

Dove serve un colore \enquote{uguale al fondo} --- l'intestazione fissa della
tabella, la maniglia delle metriche e la sua sfumatura, i cerchi dei passi,
l'anello dei punti del mirino --- si usa \cod{--ovx-bg}, cioè
\cod{--qs-canvas}; il passaggio del mouse è \cod{--ovx-hover}
(\cod{--ui-panel} sul fondo, \cod{--ui-panel-hover} dentro menu e pop-up, che
tengono il colore dei pannelli). Le metriche stanno in una griglia con il
bordo: le linee interne le disegna ogni casella con un'ombra a destra e sotto
(\cod{--ovx-cell}), e quelle dell'ultima colonna e dell'ultima riga finiscono
sotto il bordo della griglia, che le copre. Se l'ultima riga non è piena, lo
spazio vuoto resta del colore della pagina.

La colonna di destra è \cod{clamp(400px, 31\%, 468px)}. Sopra i \num{1100}
pixel di larghezza e i \num{640} di altezza il canvas della pagina prende
\cod{calc(100vh - var(--qs-tb-h))}, come le altre pagine a tutta altezza
(\vediSez{studio-pagine}), senza margine in basso: le sezioni arrivano al bordo
della finestra e la pagina non scorre. Sotto, le colonne diventano
\cod{display: contents} e le sezioni si dispongono con
\cod{grid-template-areas}: una sotto l'altra, o su due colonne fra \num{760} e
\num{1100} pixel, con un \cod{border-top} ciascuna (e un \cod{border-left} sulle
schede quando stanno accanto alle metriche). In questa disposizione la tabella
non scorre per conto suo, perché senza un pannello attorno non si vedrebbe
dove finisce: scorre la pagina, e l'intestazione resta ferma sotto la striscia
in alto (\cod{top: var(--qs-tb-h)}). Le colonne della tabella lasciano il
posto con le container query (\cod{@container ovstreams}): via la
correlazione sotto \num{860} pixel, la curva sotto \num{760}, i fold sotto
\num{640}, il drawdown sotto \num{520}.

\subsection{Regole del foglio di stile}

Ogni regola di \file{home.css} parte da \cod{\#ov-root}, tranne quelle sul
canvas della pagina (\cod{\#qs-canvas[data-tab="home"]}): lo controlla
\file{prova\_overview.py}. È la stessa ragione delle altre pagine rifatte ---
\file{ui-system.css} arriva \emph{dopo} il foglio della pagina --- con tre
conseguenze pratiche:

\begin{description}[leftmargin = 1.4em, style = nextline, itemsep = 4pt]
  \item[\cod{[hidden]}] \cod{\#ov-root [hidden]} è \cod{display: none} con
    \cod{!important}: una regola come \cod{.ovx-btn \{ display: inline-flex \}}
    altrimenti mostrerebbe un pulsante nascosto.
  \item[Selezione unica a contrasto pieno] Periodo, schede, scelte dei filtri:
    il selezionato ha \cod{--ui-solid-bg}/\cod{--ui-solid-tx}, come nel resto
    dell'app. Le schede sono \cod{role="tablist"}, e la regola globale su
    \cod{[aria-selected="true"]} dà già quel fondo.
  \item[Pulsanti] I contenitori portano \cod{data-ui-group="none"} e le
    regole dei pulsanti hanno \cod{!important} dove \file{ui-system.css} ne
    usa uno (raggio, fondo dei premuti \cod{aria-pressed}). Con l'id davanti
    vincono comunque.
\end{description}

% La tabella delle fonti non si spezza: se non sta in fondo alla pagina, meglio
% che la sottosezione cominci su quella dopo (lasciando bianco in fondo) che
% una pagina con gli spazi allargati per arrivare al margine.
\filbreak
\subsection{Lo stato e le fonti}

\file{home.js} tiene uno stato \cod{S} e una funzione di disegno per ogni
parte; nessuna parte legge il DOM di un'altra.

\begin{center}
\begin{tabular}{@{}L{46mm}L{92mm}@{}}
\toprule
\textbf{Fonte} & \textbf{Che cosa ne esce} \\
\midrule
\rigaTenue
\percorsoHTTP{/api/a/overview?spark=1} & Le righe della tabella: Sharpe, CAGR,
  drawdown, correlazione con SPY, fold, PSR, e la curva in 64 punti (solo
  con \cod{spark=1}); i passi della guida; il target di volatilità. \\
\percorsoHTTP{/api/projects} & I progetti, con le metriche salvate da
  Ensemble. \\
\rigaTenue
\percorsoHTTP{/api/ensemble\_custom} & Per il progetto mostrato:
  \cod{equity\_curve}, \cod{spy\_equity} e \cod{drawdown} sulle stesse date,
  \cod{metrics}, \cod{robustness}, \cod{yearly}, \cod{stream\_weights},
  \cod{weights\_info}, \cod{correlation\_matrix},
  \cod{drawdown\_attribution.\_summary}. Le risposte restano in
  \cod{S.ensCache}. \\
\percorsoHTTP{/api/overview/prefs} & Progetto scelto e metriche visibili, in
  ordine. \\
\rigaTenue
\percorsoHTTP{/api/metrics}, \cod{/compute}, \percorsoHTTP{/api/a/metric\_preview}
  & Le metriche personalizzate, i loro valori sul progetto, l'anteprima nel
  costruttore. \\
\bottomrule
\end{tabular}
\end{center}

\noindent
La rotta passa allo script anche i nomi dei modi di pesatura
(\cod{data-modes}, da \simbolo{EnsembleEngine.weighting\_modes\_meta}): la
pagina scrive \enquote{Inverse-Cov}, non \cod{inverse\_cov}. Le preferenze
di comodità stanno in \cod{localStorage}: \cod{lyra-ov-chart} (periodo, SPY,
drawdown), \cod{lyra-ov-filters} (ordinamento e filtri), \cod{lira\_fav\_streams}
(le stelle), \cod{lyra-ov-side} (la scheda aperta).

\subsection{Quale progetto: \enquote{Automatic}}

Senza una scelta salvata la pagina mostra il progetto con lo Sharpe più alto
(\simbolo{pickProject}). Lo Sharpe è quello calcolato se il progetto è già in
\cod{S.ensCache}, altrimenti quello salvato da Ensemble. Un progetto vecchio,
salvato senza metriche, sarebbe escluso dal confronto senza che nessuno lo
sapesse: \simbolo{automatico} calcola prima i progetti di cui lo Sharpe non si
conosce (al più otto, in parallelo), poi sceglie. Con le metriche salvate,
il caso normale, non parte nessuna chiamata in più.

\subsection{Il grafico}

Un SVG la cui \cod{viewBox} è la sua misura in pixel, ridisegnato da un
\cod{ResizeObserver}: testi e tratti non si deformano, e i colori vengono dai
token del tema, quindi il cambio di tema non chiede di ridisegnare.
\simbolo{windowData} taglia la finestra del periodo (22, 126, 252, 756 sedute
o tutto) e allinea SPY e drawdown per posizione, o per data se le lunghezze
non coincidono. La scala di destra è la variazione dalla partenza della
finestra, con un passo \enquote{tondo} (\simbolo{passo}) scelto per avere una
tacca ogni 34 pixel circa; il drawdown è sotto, in una fascia sua, e resta
quello di tutta la storia (la distanza dal massimo vero, non da quello della
finestra).

Il mirino non ridisegna: sposta le linee e le etichette che il disegno ha già
creato, e nasconde le tacche e le etichette fisse che finirebbero sotto le
sue. Il numero grande ha tre modi (\simbolo{readout}): \cod{scrub} segue il
mirino, \cod{rest} torna all'ultimo valore senza animazione, \cod{anim} ci
arriva contando (cambio di progetto). Le frecce, \texttt{Home}, \texttt{End} ed
\texttt{Esc} fanno lo stesso da tastiera, con l'SVG nel giro del fuoco.

\subsection{Menu e finestre}

I menu (progetto mostrato, progetto della tabella, filtri) sono
\cod{role="menu"}. All'apertura il fuoco va sulla voce scelta, le frecce
scorrono le voci, \texttt{Esc} chiude e riporta il fuoco sul pulsante. Il
verso si decide all'apertura (\simbolo{placeMenu}): verso l'alto se sotto non
c'è posto per il menu più lungo che può diventare, e poi non cambia --- un
menu che salta dall'altra parte mentre ci si clicca dentro fa perdere il
segno. Le due finestre (metriche e idoneità) tengono il fuoco dentro finché
sono aperte.

\subsection{Le prove}

\file{tests/studio\_finto.py} con \cod{--stream} prepara otto strategie con
i risultati fuori campione dal 2019 (con il benchmark e un fattore comune,
così un progetto ha uno Sharpe credibile), i fold del walk-forward per
cinque, tre progetti con modi di pesatura veri e i primi due passi della
guida fatti dall'utente della sessione. I nomi evitano le parole dei pattern
degli stream storici di \file{dashboard/config.py}: un file
\file{wf\_crypto\_momentum\_*} sarebbe diventato anche lo stream
\enquote{momentum}.

\file{tests/prova\_overview.py} controlla senza browser i dati di prova, la
pagina (ogni id cercato dallo script esiste), le API con i campi che la
pagina legge, le preferenze e il foglio di stile. \file{tests/e2e\_overview.js}
avvia due copie dell'app, una con i dati e una vuota, e prova in Chrome la
disposizione nei due temi (nessun contenitore né fondo proprio, il fondo
interno uguale a quello della pagina, le linee, gli allineamenti al titolo e
ai bordi), il portafoglio (menu, periodi, serie, mirino col mouse
e con la tastiera), le metriche (valori uguali all'API, riordino, una metrica
creata ed eliminata), le schede, la tabella (ordinamento, ricerca, filtri,
progetti, stelle, idoneità) e le finestre strette (sezioni impilate, tabella
che scorre con la pagina, intestazione ferma sotto la striscia).

\begin{codiceInterno}{console}{Le prove della Overview}{}
python tests/prova_overview.py
PLAYWRIGHT_DIR=... CHROME_BIN=/usr/bin/google-chrome node tests/e2e_overview.js
python tests/studio_finto.py --porta 5075 --sessione /tmp/sessione.txt --stream
\end{codiceInterno}

\noindent
L'ultima riga lascia accesa l'app con gli stream e i progetti, per lavorare
sull'aspetto della pagina; l'indirizzo con il gettone si stampa all'avvio e
porta direttamente alla Overview.
